\section{Experiments}
\label{sec:experiments}

\subsection{Dataset and Splits}

The CQ500 dataset~\cite{chilamkurthy2018cq500} contains 491 non-contrast
head CT studies contributed from emergency departments.
Table~\ref{tab:dataset} summarises the study-level class distribution after
majority-vote label assignment.
Epidural hemorrhage is the rarest subtype, which creates significant class
imbalance that is addressed through positive class weighting
(Section~\ref{sec:loss}).
Studies are partitioned at the study level in a stratified $70/15/15$ split
using \emph{any hemorrhage} as the stratification variable, ensuring that the
positive rate is consistent across splits.
No data augmentation is applied to the validation or test splits.

\begin{table}[t]
    \centering
    \caption{CQ500 study-level class distribution. Percentages are relative
             to the 491 total studies.}
    \label{tab:dataset}
    \begin{tabular}{lcc}
        \toprule
        \textbf{Class} & \textbf{Positive studies} & \textbf{\%} \\
        \midrule
        Any hemorrhage      & 205 & 41.8 \\
        Intraparenchymal    & 102 & 20.8 \\
        Subdural            & 117 & 23.8 \\
        Subarachnoid        &  73 & 14.9 \\
        Intraventricular    &  62 & 12.6 \\
        Epidural            &  25 &  5.1 \\
        \bottomrule
    \end{tabular}
\end{table}

\subsection{Experimental Conditions}
\label{sec:conditions}

We evaluate the following systems:

\begin{enumerate}
    \item \textbf{Baseline -- DenseNet-121 only}: a single DenseNet-121
          backbone trained with the adjacent-slice 3-channel input, no
          sequence model.
    \item \textbf{Ensemble (Stage 1 only)}: average of DenseNet-121,
          DenseNet-169, and SE-ResNeXt-101 predictions, without Stage-2.
    \item \textbf{Ensemble + Sequence (Stage 1 + 2, without $\delta$)}:
          full two-stage pipeline without slice thickness input.
    \item \textbf{Ensemble + Sequence + $\delta$ (full model)}: full
          two-stage pipeline including the slice thickness feature in SM2
          (Eq.~\ref{eq:sm2_input}).
\end{enumerate}

Systems~3 and~4 correspond to the original first-place RSNA 2019 solution
without and with the slice thickness feature, respectively.
All systems share the same preprocessing pipeline, loss function, and
optimiser configuration described in Section~\ref{sec:training}.

\subsection{Hardware and Software}

All experiments are run on a single workstation with one NVIDIA GPU
(minimum 8\,GB VRAM) under Ubuntu 22.04.
The software stack is: Python 3.11, PyTorch 2.x, torchvision, timm,
albumentations, scikit-learn, pydicom, and MLflow~\cite{mlflow2018} for
experiment tracking.
Reproducibility is ensured via a DVC~\cite{ruslan2020dvc} pipeline that
records data preprocessing, model training, and evaluation as separate
tracked stages with content-hashed outputs.

\subsection{Evaluation Metrics}
\label{sec:metrics}

We report the following metrics on the held-out test set.

\paragraph{Area under the ROC curve (AUC).}
Per-class AUC is computed against continuous probability predictions.
Macro-AUC is the unweighted mean over all $C = 6$ classes.

\paragraph{Sensitivity and Specificity.}
At a decision threshold $\tau = 0.5$, we compute per-class sensitivity
(true positive rate) and specificity (true negative rate):
\begin{align}
    \mathrm{Sensitivity}_c &= \frac{\mathrm{TP}_c}{\mathrm{TP}_c + \mathrm{FN}_c}, \nonumber \\
    \mathrm{Specificity}_c &= \frac{\mathrm{TN}_c}{\mathrm{TN}_c + \mathrm{FP}_c}.
    \label{eq:sens_spec}
\end{align}

\paragraph{Bootstrap 95\,\% Confidence Intervals.}
For each AUC value, we apply percentile bootstrap resampling with
$N_{\mathrm{boot}} = 1000$ resamples, drawing with replacement from the
test-set slice predictions.
The 2.5th and 97.5th percentiles of the bootstrap AUC distribution form
the 95\,\% CI~\cite{efron1994introduction}.
This follows the reporting methodology of Wang et al.~\cite{wang2021deeplearning}.
